\documentclass[a4paper]{article}
\usepackage{amsfonts}
\usepackage{amsmath}
\usepackage{amssymb}
\usepackage{graphicx}     

% Command definities van frequent gebruikte commandos
\newcommand{\N}{\mathbb{N}}
\newcommand{\Q}{\mathbb{Q}}
\newcommand{\R}{\mathbb{R}}
\newcommand{\C}{\mathbb{C}}
\newcommand{\Bgsin}{\operatorname{Bgsin}}
\newcommand{\Bgcos}{\operatorname{Bgcos}}
\newcommand{\Bgtan}{\operatorname{Bgtan}}
\newcommand{\cis}{\operatorname{cis}}

\begin{document}
  
\noindent \large Auteur: Wout Vanderborght\\
\noindent \large Studierichting: Informatica\\
\noindent \large Oefeningenreeks 2A nr. 2.2\\

\medskip

\normalsize

Bepaal van de volgende functie $f: \R \rightarrow \R$ het domein: \\

$ f(x) = \dfrac{5x}{2x + 3} $ als $ x < 1 $ \\

Nulpunt noemer: $ 2x + 3 = 0 \Leftrightarrow x = -\dfrac{3}{2} $ \\

$ \operatorname{dom} f = \left]-\infty, 1\right[ \setminus \left\{-\dfrac{3}{2}\right\} $ \\

\begin{thebibliography}{999}
%\bibitem{a} Referentie
\end{thebibliography}
\end{document} 
